El presente trabajo no pretende cubrir todo el espacio de modelos descrito en el marco teórico, sino seleccionar el subconjunto mínimo suficiente para producir resultados medibles y reproducibles sobre un fenómeno de guerra cognitiva. Se combinan dos topologías contrastadas con el agente de umbral lineal.

La \textbf{red libre de escala} (\emph{scale-free}, variante digital dirigida) modela el entorno de las plataformas digitales. Su distribución de grado sigue una ley de potencias con exponente $\gamma \in [2, 3]$ \cite{BarabasiAlbert1999}, concentrando la influencia en un pequeño número de \emph{hubs} que reproducen la asimetría seguidor-seguido característica de Twitter o Telegram. Esta topología constituye el sustrato donde la guerra cognitiva opera con mayor eficacia documentada \cite{Vosoughi2018}.

La \textbf{red de mundo pequeño} (\emph{small-world}, variante analógica bidireccional) modela el tejido de relaciones presenciales y de confianza cercana, generada mediante el modelo de Watts-Strogatz \cite{WattsStrogatz1998} con $k=6$ vecinos y probabilidad de recableado $p=0{,}10$. Su alto coeficiente de agrupamiento local, sin \emph{hubs} de alto grado, representa la capa analógica de la sociedad donde la difusión es más lenta pero el refuerzo social entre pares cercanos es más intenso.

Incluir ambas topologías en el diseño experimental permite contrastar cómo la estructura de la red modula la resistencia cognitiva: la red libre de escala maximiza el alcance del ataque vía \emph{hubs}, mientras que la red de mundo pequeño ralentiza la propagación a través del efecto de agrupamiento local.

De los distintos arquetipos de agente, se elige el \emph{agente resistente}, que incorpora un umbral explícito de adopción $\theta_i$ y acumula convicción $c_i(t)$ de forma determinista según la ecuación \ref{eq:conviction}. Este modelo es el mínimo con cognición medible: proporciona variables de estado internas interpretables, no requiere calibración estocástica adicional más allá del umbral, y produce cascadas cuya estructura puede analizarse sin ambigüedad.

\medskip

El problema se plantea formalmente de la siguiente manera. Sea $G = (V, E)$ una red libre de escala dirigida con $|V| = N$ nodos. En el instante $t = 0$ se inyecta una narrativa a través de un conjunto de semillas $S \subset V$, con $|S| = k$, que quedan con convicción unitaria y son consideradas adoptantes. El resto de la población parte de $c_i(0) = 0$. A lo largo de la simulación discreta se monitoriza el conjunto de adoptantes $A(t) = \{ i \in V \setminus S : c_i(t) \geq \theta_i \}$, donde $c_i(t)$ es la convicción del agente $i$ en el paso $t$ tal como se define en \ref{eq:conviction}. El objetivo del experimento es caracterizar la dinámica de $|A(t)|$ y cuantificar la resistencia cognitiva de la población ante distintas condiciones de ataque.

\medskip

La heterogeneidad en el umbral de adopción es el principal determinante estructural de si una cascada se detiene o se autopropaga \cite{Granovetter1978, Watts2002}. Para reflejar la coexistencia de dos subpoblaciones ---susceptibles e informacionalmente resistentes--- el umbral de cada agente se muestrea de una distribución bimodal según la ecuación

\begin{equation}
  \label{eq:theta}
  \theta_i \sim
  \begin{cases}
    \operatorname{clip}\!\left(\mathcal{N}(\mu_s,\,\sigma^2),\,0,\,1\right) & \text{con probabilidad } p_s,\\[4pt]
    \operatorname{clip}\!\left(\mathcal{N}(\mu_r,\,\sigma^2),\,0,\,1\right) & \text{con probabilidad } 1 - p_s,
  \end{cases}
\end{equation}

\noindent donde $\mu_s < \mu_r$, de modo que los agentes susceptibles adoptan la narrativa ante exposiciones de menor intensidad que los resistentes, y la función $\operatorname{clip}(\cdot, 0, 1)$ garantiza que el umbral permanece en el dominio válido. Los valores de referencia son $\mu_s = 0{,}30$, $\mu_r = 0{,}70$ y $\sigma = 0{,}10$.

\medskip

La narrativa es inyectada por las semillas en $t = 0$ mediante un mecanismo de \emph{fanout parcial}: cada semilla selecciona $\lfloor \rho \cdot |\text{seguidores}| \rfloor$ destinos de entre sus vecinos salientes, donde $\rho \in (0, 1]$ es la fracción de audiencia contactada en el primer paso. Este diseño evita que la semilla alcance a toda su red de golpe, reproduciendo el comportamiento real de la publicación inicial en redes sociales, donde la difusión inicial afecta a una fracción de la audiencia total antes de activar mecanismos de amplificación orgánica. Se contemplan dos estrategias de selección de semillas: la estrategia \emph{hubs}, que elige los $k$ nodos de mayor in-grado, y la estrategia \emph{aleatoria}, que los escoge uniformemente al azar; el contraste entre ambas operacionaliza el problema del sembrado óptimo estudiado en \cite{KempeKleinbergTardos2003}.

\medskip

Para dotar al modelo de un periodo infeccioso superior a un único paso de simulación ---condición necesaria para que los adoptantes contribuyan activamente a la propagación--- se incorpora la \emph{re-emisión espontánea}: en cada paso $t$, todo agente que haya adoptado la narrativa la re-emite con probabilidad $f$ hacia un subconjunto de sus vecinos salientes. Este mecanismo es coherente con la teoría del contagio complejo \cite{CentolaMacy2007}, según la cual la difusión de creencias requiere exposiciones reiteradas a través de múltiples contactos, y modela el comportamiento observado de compartición recurrente de contenido en plataformas digitales.

\medskip

La resistencia cognitiva de la población se cuantifica mediante cuatro métricas complementarias. La \emph{tasa de ataque} $a = |A(T)| / (N - k)$, donde $T$ es el instante final de la simulación, mide la fracción de la población no semilla que ha sido convertida; la \emph{resiliencia poblacional} se define como $R = 1 - a$. La \emph{entropía cognitiva} en bits,

\begin{equation}
  \label{eq:entropia}
  H = -\bigl[a \log_2 a + (1 - a) \log_2 (1 - a)\bigr],
\end{equation}

\noindent alcanza su máximo ($H = 1$) cuando la población está dividida en partes iguales y tiende a cero cuando el consenso colapsa hacia uno de los dos estados, reflejando la pérdida de pluralidad cognitiva en el sentido del Cognitive Net Assessment \cite{PolytechniqueCNA2026}. A nivel individual, la \emph{resistencia individual} del agente $i$ se define como

\begin{equation}
  \label{eq:Ri}
  R_i =
  \begin{cases}
    1 & \text{si el agente no adopta nunca,}\\[4pt]
    1 - \dfrac{1}{d_i} & \text{si adopta tras } d_i \text{ exposiciones,}
  \end{cases}
\end{equation}

\noindent donde $d_i$ es la dosis acumulada en el momento de la conversión (no la acumulada al final de la simulación), de modo que un agente que cede en la primera exposición obtiene $R_i = 0$ y uno que resiste muchas exposiciones antes de ceder obtiene $R_i$ próximo a la unidad. Por último, el \emph{tiempo de penetración} $t_{50}$ es el paso en que se ha convertido la mitad de los agentes que acabarán convirtiéndose, y la curva $|A(t)|$ registra la trayectoria completa de adopción.

\medskip

Los hiperparámetros del experimento se fijan según los valores de la Tabla~\ref{tab:hiperparametros}, que recoge su valor de referencia, la justificación funcional y la fuente bibliográfica que respalda la elección.

\begin{table}[ht]
  \centering
  \small
  \caption{Hiperparámetros del experimento de referencia.}
  \label{tab:hiperparametros}
  \begin{tabular}{@{}p{3.2cm} p{2.4cm} p{5.2cm} l@{}}
    \toprule
    \textbf{Hiperparámetro} & \textbf{Valor} & \textbf{Justificación} & \textbf{Fuente} \\
    \midrule
    $\mu_s$, $\mu_r$, $\sigma$ & $0{,}30$ / $0{,}70$ / $0{,}10$ & Dos subpoblaciones; la heterogeneidad gobierna las cascadas & \cite{Granovetter1978,Watts2002} \\
    \addlinespace
    $p_s$ (fracción susceptible) & $0{,}50$ (barrida) & Composición de la población; variable de la curva dosis-respuesta & \cite{Watts2002} \\
    \addlinespace
    Ganancia $\gamma$ & $0{,}50$ & Refuerzo gradual por exposición (contagio complejo) & \cite{CentolaMacy2007} \\
    \addlinespace
    Carga emocional $\lambda$ & $0{,}80$ & Único \emph{driver} de persuasión; alta activación difunde más & \cite{Vosoughi2018} \\
    \addlinespace
    Umbral de veracidad $v_{\mathrm{thr}}$ & $0{,}30$ & Define la desinformación rastreada & \cite{Vosoughi2018} \\
    \addlinespace
    Prob.\ de reenvío $f$ & $0{,}10$ & Periodo infeccioso $> 1$ paso sin saturar la red & \cite{KempeKleinbergTardos2003} \\
    \addlinespace
    Prob.\ de ruido $p_{\mathrm{fire}}$ & $0{,}05$ & Canal ambiental vivo & --- \\
    \addlinespace
    Estrategia de semilla & \emph{hubs} / aleat. & Problema de \emph{seeding} & \cite{KempeKleinbergTardos2003} \\
    \addlinespace
    $k$ semillas / fanout $\rho$ & $25$ / $0{,}4$ & Ataque concentrado con difusión parcial & --- \\
    \addlinespace
    Libre de escala $\alpha,\beta,\gamma_{\mathrm{BA}},\delta_{\mathrm{in}}$ & $0{,}41$/$0{,}54$/$0{,}05$/$0{,}2$ & Producen $\gamma \in [2, 3]$ típico de redes digitales dirigidas & \cite{Bollobas2003} \\
    \addlinespace
    Mundo pequeño $k$, $p_{\mathrm{rewire}}$ & $6$ / $0{,}10$ & Alto agrupamiento y diámetro corto; capa analógica & \cite{WattsStrogatz1998} \\
    \addlinespace
    $N$ / réplicas & $5\,000$ / $25$ & Compromiso coste-realismo; Monte Carlo sobre el \emph{ensemble} & --- \\
    \bottomrule
  \end{tabular}
\end{table}

\medskip

El diseño experimental sigue la estrategia \emph{one-factor-at-a-time} (OFAT) combinada con Monte Carlo sobre el \emph{ensemble}: se parte de una configuración de referencia (\emph{baseline}) y se varía un único factor por experimento, manteniendo el resto fijo. Cada celda experimental se repite con $25$ semillas pseudoaleatorias distintas; cada semilla regenera de forma independiente el grafo, los umbrales $\theta_i$ y la trayectoria dinámica, produciendo muestras estadísticamente independientes. Los resultados se reportan como media $\pm$ desviación estándar sobre las réplicas, nunca como una corrida única. El barrido cubre cinco factores de la narrativa sobre las dos topologías, lo que produce $2 \times (3+2+3+3+3) \times 25 = 700$ runs totales ($\approx 30$ minutos de CPU).

\begin{table}[ht]
  \centering
  \small
  \caption{Diseño OFAT del barrido experimental (sweep v2). Cada celda se repite con 25 réplicas Monte Carlo sobre las dos topologías.}
  \label{tab:experimentos}
  \begin{tabular}{@{}l p{2.6cm} p{4.5cm} c p{3.4cm}@{}}
    \toprule
    \textbf{Factor} & \textbf{Símbolo} & \textbf{Niveles (baseline en negrita)} & \textbf{Niv.} & \textbf{Métricas registradas} \\
    \midrule
    Ganancia persuasión & $\gamma$ & $0{,}25$; \textbf{$0{,}50$}; $0{,}75$ & 3 & $a$, $R$, $H$, $\bar{R}_i$, $t_{50}$, $|A(t)|$ \\
    \addlinespace
    Estrategia semilla & --- & \textbf{\emph{hubs}}; aleatorio & 2 & $a$, $R$, $H$, $\bar{R}_i$, $t_{50}$ \\
    \addlinespace
    Fanout inicial & $\rho$ & $0{,}2$; \textbf{$0{,}4$}; $0{,}6$ & 3 & $a$, $R$, $H$, $\bar{R}_i$, $t_{50}$ \\
    \addlinespace
    Número de semillas & $k$ & $5n$ $(0{,}1\%)$; \textbf{$25n$} $(0{,}5\%)$; $50n$ $(1\%)$ & 3 & $a$, $R$, $H$, $\bar{R}_i$, $t_{50}$ \\
    \addlinespace
    Prob.\ reenvío & $f$ & $0{,}10$; \textbf{$0{,}25$}; $0{,}40$ & 3 & $a$, $R$, $H$, $\bar{R}_i$, $t_{50}$ \\
    \bottomrule
  \end{tabular}
\end{table}
